\documentclass[10pt,letterpaper]{article}

\usepackage[margin=0.5in]{geometry}
\usepackage[T1]{fontenc}
\usepackage{lmodern}
\usepackage{xcolor}
\usepackage{enumitem}
\usepackage{tabularx}
\usepackage[hidelinks]{hyperref}
\usepackage{titlesec}
\usepackage{microtype}

\definecolor{accent}{HTML}{145A78}
\setlength{\parindent}{0pt}
\setlength{\parskip}{0pt}
\setlist[itemize]{leftmargin=1.15em,itemsep=0.8pt,topsep=1pt,parsep=0pt}
\titleformat{\section}{\large\bfseries\color{accent}}{}{0pt}{}[\vspace{-2pt}\titlerule]
\titlespacing*{\section}{0pt}{5pt}{2pt}
\pagestyle{empty}

\newcommand{\role}[4]{%
  \textbf{#1}\hfill #2\\
  \textit{#3}\hfill \textit{#4}\vspace{-3pt}
}

\begin{document}
\fontsize{9}{10.5}\selectfont

\begin{center}
  {\Large\bfseries Piter Garcia}\\[1pt]
  {\large Elementary Inclusive Educator and Computer Science Teacher Candidate}\\[2pt]
  Rochester, NY \enspace|\enspace Available for NYC relocation beginning January 2027\\
  \href{mailto:pgarcia8@u.rochester.edu}{pgarcia8@u.rochester.edu}
  \enspace|\enspace +1 (646) 288-3300
  \enspace|\enspace \href{https://linkedin.com/in/garciapiterz}{linkedin.com/in/garciapiterz}\\
  \href{https://github.com/pzg8794}{github.com/pzg8794}
  \enspace|\enspace \href{https://garciapiterz.com}{garciapiterz.com}
\end{center}

\section*{Profile}
K--5 educator and NSF Robert Noyce Scholar completing an M.S. in Teaching
Computer Science K--12 at the University of Rochester. Brings supervised
student-teaching experience across kindergarten through fifth grade together
with graduate preparation in disability and inclusive practices, Universal
Design for Learning, differentiation, MTSS, IEP-informed instruction,
assessment, behavior supports, and culturally responsive teaching. Prior
software, data, and AI experience adds a disciplined approach to formative
evidence, progress monitoring, problem solving, and clear instructional
systems.

\section*{Elementary Teaching and Inclusive Practice}
\role{Computer Science Student Teacher}{Spring 2026}
{University of Rochester / Pine Brook Elementary School}{Rochester, NY}
\begin{itemize}
  \item Supported and increasingly led instruction across kindergarten through
    fifth grade using coding, robotics, Sphero, Minecraft Education, Tinkercad,
    Code.org, circuits, digital fluency, and age-appropriate AI activities.
  \item Designed and revised lessons, directions, rubrics, and formative
    assessments using Universal Design for Learning and multiple pathways for
    students to access content and demonstrate understanding.
  \item Tutored individuals and small groups, co-taught lesson segments, and
    adjusted pacing, grouping, visuals, hands-on supports, and task structure in
    response to student readiness and classroom evidence.
  \item Used clear routines, relationship-building, engaging entry points, and
    structured choices to help students move from off-task behavior toward
    sustained work, successful participation, and peer support.
  \item Worked within a supervised placement that included both computer-science
    and inclusion/special-education guidance, with increasing responsibility
    for planning, instruction, assessment, and classroom management.
\end{itemize}

\role{Tutor and Mentor}{Ongoing}
{Computer Science, Data Science, Statistics, and Mathematics}{Rochester, NY}
\begin{itemize}
  \item Break complex tasks into manageable steps using examples, diagrams,
    guided practice, and patient debugging while preserving learner ownership.
  \item Support neurodivergent and underrepresented learners with organization,
    feedback interpretation, confidence-building, and flexible ways to engage
    with rigorous technical work.
\end{itemize}

\section*{Inclusive-Education Preparation}
\begin{tabularx}{\textwidth}{@{}>{\bfseries}l X@{}}
Disability and inclusion & Graduate study in disability and schools, inclusive classroom practice, teaching students with disabilities, and collaboration across educational roles\\
Instructional design & Universal Design for Learning, differentiation, MTSS, scaffolding, learner agency, and culturally responsive pedagogy\\
IEP-informed practice & Coursework and supervised preparation involving IEP goals, accommodations and modifications, assessment adaptation, progress monitoring, and team-based decision making\\
Behavior and regulation & Proactive routines, self-regulation supports, reflection, classroom climate, restorative approaches, and behavior as instructional information\\
Assessment & Formative evidence, varied response modes, progress checks, and instructional adjustment based on observed student needs\\
\end{tabularx}

\section*{Research and Technical Experience}
\role{Graduate Assistant, Quantum and AI Research}{Fall 2025--Present}
{Rochester Institute of Technology, Software Engineering}{Rochester, NY}
\begin{itemize}
  \item Build and document reproducible Python experiments in adaptive
    decision-making, fairness, and resource allocation; translate complex
    methods into clear reports and reusable learning materials.
  \item Apply structured validation, logging, comparison, and reflection cycles
    that transfer naturally to data-informed instructional planning.
\end{itemize}

\role{Data Solutions Engineer}{Sep 2022--Aug 2024}
{VEDADATA Inc.}{New York, NY}
\begin{itemize}
  \item Built maintainable Python data workflows, validation checks, and
    documentation while collaborating across technical and nontechnical teams.
\end{itemize}

\role{AI Data Solutions Engineer}{Jan 2021--Sep 2022}
{VIOME Inc.}{New York, NY}
\begin{itemize}
  \item Developed and evaluated machine-learning workflows for healthcare data,
    strengthening skills in careful measurement, iteration, and communication.
\end{itemize}

\section*{Education}
\textbf{M.S., Teaching Computer Science K--12}, University of Rochester
\hfill Expected December 2026\\
NSF Robert Noyce Scholar; graduate preparation in disability and inclusive practices

\textbf{M.S., Data Science}, Rochester Institute of Technology
\hfill Expected December 2026\\
Applied AI, bioinformatics, neural networks, statistics, and reproducible data systems

\textbf{M.S., Computer Science}, Rochester Institute of Technology
\hfill 2015\\
Artificial intelligence, big-data analytics, and computer graphics

\textbf{B.S., Electrical Engineering Technology; B.S., Computer Engineering Technology},
Farmingdale State College \hfill 2012

\section*{Teaching Strengths}
Elementary instruction \enspace|\enspace UDL and differentiation \enspace|\enspace
IEP-informed supports \enspace|\enspace Formative assessment \enspace|\enspace
Small-group instruction \enspace|\enspace Behavior and self-regulation supports\\
Culturally responsive teaching \enspace|\enspace Family/team collaboration \enspace|\enspace
Computer science and STEM enrichment \enspace|\enspace Data-informed reflection

\end{document}
